\chapter{Literature Review}\label{ch:Background}

\section{Chapter Overview}

This chapter aims at critically discussing the scholarly literature on AI-based supply chain analytics in the context of consumer goods, with the aim of giving a holistic view of theoretical basis, technological inventions, and research outcomes upon which this study draws upon. Integrating studies on the use of machine learning, as well as hybrid intelligence, and the use of data analytics in supply chain management, this review proposes the novelty of research gaps in the literature, emphasizing methods, and locating the study in the context of the larger academic discussion of digital transformation in logistics. The three thematic dimensions (technological innovation, organizational transformation, and performance outcomes) used to organize the chapter enables to have a clear idea of how the AI-human-collaboration can contribute to improving supply chain visibility, flexibility, and sustainability. By means of this synthesis, the review provides the theoretical and practical premises to the development of a hybrid AI-based analytics tool, in which the needs of global consumer goods supply chains will be adjusted to the new requirements, and key gaps in the concept of ethical AI implementation, data integration challenges, and emerging market uses, which can justify the necessity to conduct additional research, are identified.


\section{Background to the Problem}

Consumer goods companies such as Unilever and Nestle supply chain management are being increasingly focused on proper demand forecasting and nimble decision-making to control inventory and costs \autocite{luo_intelligent_2024}. The conventional time-series models, like ARIMA and Holt-Winters, do not perform well when the demand is volatile and unpredictable disruptions occur. The most recent developments in artificial intelligence (AI) and machine learning (ML) such as LSTM and XGBoost have increased the accuracy of the forecast by 15 to 20 percent, and this is now applicable to inventory turnover, stockouts, and cost reduction \autocite{chowdhury_systematic_2025}. Nevertheless, there are still difficulties: AI models are generally not transparent enough to give the planners confidence, and most of them are backward-looking, which means a significant dependence on past sales and a lack of support to real-time drivers of the market, including macroeconomic or social indicators. 

Also, such standard measures as MAPE do not take into account asymmetric business costs of over- and under-forecasting. Forecasting is also not integrated with prescriptive decision systems and does not consider reverse logistics which is a key aspect in e-commerce and sustainability. To overcome these shortcomings, this proposal presents a hybrid, interpretable, and cost-sensitive model of analytics using statistical, signal-processing, and ML methods. It involves external information in real time, explainability procedures, asymmetric maximization of business goals and reverse-flow predictions. It aims to offer strong, transparent and practical forecasts that complement efficiencies, cost-effectiveness, and sustainability in consumer goods supply chains.


\section{Theoretical and Conceptual Background}

The combination of Data Science and Artificial Intelligence (AI) in supply chain analytics has altered the old, linear systems of operation to smart, self learning and adaptable networks. The AI provides predictive, informative, and strategic capabilities in the consumer goods industry, where the demand is volatile, the product life cycle is short, and logistics are complex. This section will create the theoretical base and conceptual knowledge in which the study will be based. It describes the role of AI, big data, and hybrid analytics in creating more resilient, smart, and sustainable supply chains, as well as specifying key terms, ideas, and models used to formulate the intelligent supply chain in the modern context.

Theoretically speaking, a number of frameworks describe the ways in which AI-based analytics enhances the performance of the supply chain and competitive edge. According to the Resource-Based View (RBV) theory \autocite{barney1991}, firms can experience a sustainable advantage in case they bear resources that are unique, valuable and inimitable \autocite{Jackson01092024}. These key resources in the contemporary world are digital assets, including information infrastructure, artificial intelligence, and analysis capabilities \autocite{mailani2024}. Using AI to convert gigantic quantities of raw data into predictive and prescriptive data, consumer goods companies will have successfully turned digital capability into a strategic asset that aids decision-making and continuous improvement \autocite{groenewald2024, jauhar_explainable_2024, aghaei2025potentiallargelanguagemodels}.

This is further complemented by the Dynamic Capabilities Theory \autocite{teece_explicating_2007} which provides an explanation of how the firms adapt to the volatile environments in response to the capability to sense opportunities and seize them and reorganize internal processes. These dynamic capabilities are enhanced by AI and machine learning since it enables supply chain managers to identify the changes in the market, predict disruption, and real-time adjustment of the decisions \autocite{tothz2025}. An example is AI-based demand forecasting which assists companies to respond to seasonal changes in sales, evolving customer preferences, or delays in supply to continue with the service even when uncertain \autocite{ahn2024generativeprobabilisticplanningoptimizing, kumar2024ai}.

Also, there is the Sociotechnical Systems Theory \autocite{bostrom1977}, which puts an emphasis on harmonizing technology with humans. The effectiveness of the AI integration process is not determined only by the technology itself but also by the level at which the technology can collaborate with the human expertise. This view is core to hybrid AI-based supply chains, where AI analytics is supplemented by human judgment to deal with ambiguous, ethical or context-specific decisions. Instead of substituting human beings, hybrid systems enhance human knowledge- a kind of balancing process where technology can be applied in analyzing data, whereas human beings can monitor, process, and authenticate decisions \autocite{vann2025human}. This principle of augmented intelligence creates responsible and efficient decision-making in place of a complicated supply chain in the world.

In this context, it is essential to determine the main concepts in order to comprehend AI and Data Science. Artificial Intelligence (AI) is offered as a term to describe computer systems that can replicate human cognitive processes, including reasoning, learning, and problem-solving. In the supply chain, AI is represented by technologies, such as Machine Learning (ML), Deep Learning (DL), Natural language Processing (NLP), and Computer Vision, each of which has a specific use of analysis. Examples of machine learning are models that, through historical data, learn to make predictions and become more accurate as time progresses without being coded. Supervised ML models are used in supply chain analytics to predict product demand or pricing, whereas unsupervised clustering methods can be used to group suppliers or customer groups together because of similar behavior patterns \autocite{KHEDR2024100379}. Deep learning, which is a subdivision of ML, operates on large and unstructured data, e.g., sensor data, customer reviews, or image data, and allows identifying defects and making dynamic predictions in logistics \autocite{hosseinnia2023applications, riesen2025sociotechnological}.

Other necessary concepts are predictive analytics and prescriptive analytics. Predictive analytics involves the application of algorithms to forecast events ahead such as sales volumes, supplier delays or stockouts whereas prescriptive analytics goes a step further by suggesting the best decisions to make in different circumstances e.g. when to replenish inventory or reroute deliveries. These processes are based on the Big Data Analytics (BDA) that means gathering the information provided by various sources, such as enterprise systems, IoT sensors, and market data, into one analytical platform. Nevertheless, \autocite{kache2017challenges} observe that the quality of data, its governance, and interoperability are key to the success of AI. Even the best algorithms are unlikely to give accurate or biased results without trustworthy real-time information.

The other important concept is the concept of Hybrid Intelligence which defines the cooperation of human thought and artificial intelligence. Rather than being completely automated, hybrid systems use AI-based computation and human-based context reasoning. Hybrid intelligence in the consumer goods business will create a balance between efficiency and responsibility in the industry where ethical and sustainability factors play a role in supply decisions \autocite{vann2025human}. This practice matches the current organizational trends of focusing on transparency, human control, and ethical governance of AI implementation.

To conceptualize these notions, there are multiple frameworks that can be used to describe the adoption of AI into the supply chain system. The Supply Chain Operations Reference (SCOR) Model is one of the oldest to organize the supply chains according to the five fundamental dimensions, Plan, Source, Make, Deliver, and Return \autocite{derfoufiintegrating}. AI improves all aspects of it by using automation and analytics: predictive planning is more precise, AI-based procurement is more efficient at choosing suppliers, robotics make the manufacturing process more efficient, and real-time tracking makes the distribution process smoother. The integration of AI into the SCOR model will turn passive processes into active, data-oriented systems that learn and become adapted at all costs.

The other applicable model is the AI-Augmented Decision-Making Model that has three layers of analytics: the descriptive, predictive, and prescriptive. At the descriptive level, AI systems represent the existing performance in form of dashboards and real-time monitoring solutions. The predictive layer uses machine learning to predict disruption or demand variations whereas the prescriptive layer gives practical advice including rerouting transport or modifying production timelines \autocite{herath2023}. These suggestions are verified by managers so that decisions are not made outside of organizational context and ethics. The approach to hybrid intelligence is based on this framework, which is known as the human-in-the-loop approach.

The DIKW Hierarchy (Data–Information-Knowledge-Wisdom) is yet another description of the process of knowledge transformation in systems that rely on AI. Information derived through different sources turns into data once processed and arranged; AI can turn this into knowledge by recognizing latent patterns and correlations; human specialists can turn this knowledge into wisdom turning insights into actionable business strategies. In a hybrid supply chain, AI can speed up the up the chain of command of any supply chain, whereas the human factor provides the wisdom component with a sense of context \autocite{DEEPU2021200048}.

The new framework in digital supply chains is the Dynamic Digital Twin Model, in which AI-based simulations generate virtual images of the physical supply chain. These digital twins are always informed by the real-world data enabling managers to simulate various situations like supply disruptions or sudden demand spikes and then execute the decisions in the real-time \autocite{espinosa2024}. Digital twins have the ability to simulate the impact of alteration of a production scenario on inventory, transportation, and energy consumption in consumer goods production, and thus make decisions out of data and preemptively \autocite{barykin2020concept, abideen2021digital}.

Data-driven management also has a number of principles that govern the application of AI in supply chains. The interoperability and integration of data are building blocks because any AI system cannot work without a integrated data stream between procurement, production, logistics and sales. The clarity and explicability are also of great essence; the AI models should be readable in order to be understood by the human managers, to be validated and, in case, override automated decisions. This increases confidence and eliminates threats of algorithm bias. Scalability and real-time responsiveness are a plus in that AI systems can process large amounts of data and deliver insights in real-time, which is needed by the fast-moving consumer goods (FMCG) industry \autocite{ghodake2024enhancing}. Finally, responsible use of technology based on the ethical and sustainable frameworks of AI is necessary to be fair, inclusive, and cause minimal damage to the environment \autocite{paleti2025artificial}. The ethical AI practices enable corporate objectives in the Walmart style eco-reduction, fair labor, and green logistics.

The idea of an AI-based supply chain analytics hybrid concept could be represented in the shape of a multi-layered ecosystem by combining these theories and frameworks. The data acquisition layer is the first layer and gathers both structured and unstructured data on suppliers, production systems and markets. The second layer- AI analytics, makes use of all these datasets through machine learning to make predictive insights. The third layer -human-AI collaboration enables the managerial interpretation and AI ethical validation of the outputs. The fourth layer is a decision execution which realizes data-driven activities in supply chain functions, and the fifth layer is feedback and learning which captures performance data to re-train AI models to improve subsequent decisions. This continuous learning so ascertained is a pointer to the Dynamic Capabilities approach because it senses change, seizes opportunity and reinvents operations to ensure resistance to change and competition \autocite{herath2023}.


\section{Review of Key Methods and Technologies}

The need to implement Artificial Intelligence (AI), machine learning (ML), and advanced analytics has been predetermined by the growing number of supply chains of consumer goods and their complexity and data intensity \autocite{samuels2025examining, sharma2022role}. This section is a critical analysis of the major methods, algorithms, and tools that are used in hybrid AI-driven supply chain analytics. It talks about preprocessing methods, methods of evaluation of the model, and contrasts the merits and demerits of various algorithmic and analytical methods in the context of predictive forecasting, optimization, and decision automation.

\subsection{Data Preprocessing and Feature Engineering}

The data preprocessing is an important aspect of supply chain analytics because it helps to ensure that raw and heterogeneous data may be transformed into structured and trustworthy inputs of AI models. The enterprise systems and IoT devices, as well as social media, provide supply chains with large amounts of transactional, sensor, and unstructured textual data \autocite{sodhi2021supply}. Best preprocessing ensures the accuracy of model and avoids overfitting or bias. Data cleaning, normalization, feature extraction, and dimensionality reduction are some of the key techniques in this process \autocite{Min2022BigDataAI}.

Data cleaning in the consumer goods environment eliminates duplicate data, absent values, anomalies in order information and logistics files \autocite{guo2023normal}. Normalization makes the various measurement scales comparable whereas feature engineering isolates important attributes like lead time variability, shelf- life or supplier reliability that enhance better model learning. Principal Component Analysis (PCA), as well as t-Distributed Stochastic Neighbor Embedding (t-SNE) is commonly used to lower the dimensions and define the most significant variables with their impact on demand and delivery \autocite{Jahin_2024}.

Time-series decomposition is another preprocessing method that is required to isolate seasonal and trend elements of the sales or inventory data, and to enable predictive models to be more understandable. Moreover, within the hybrid AI, data fusion is used in order to integrate structured ERP data with real-time IoT sensor measurements (e.g., temperature, humidity, vehicle speed) and unstructured social data (consumer reviews, social trends) to increase the level of contextual intelligence \autocite{ivanov_digital_2025}. Correct preprocessing therefore forms a solid data base that enhances the efficiency of the following machine learning and optimization models.


\subsection{Machine Learning Methods}

Machine learning algorithms are being massively employed in supplier chain forecasting, demand sensing, and anomaly detection of consumer goods. Supervised models of learning are generally employed in case of labelled data, whereas unsupervised and reinforcement learning are employed in clustering and adaptive decision-making.

Linear Regression and Multiple Regression are some of the basic models that are used in identifying the relationship that exists between demand and any force that influences the demand like price, promotion or seasonality. These models however are linear and independent of variables and therefore, they are less predictive to complex nonlinear situations. High-level algorithms like Random Forests (RF) and Gradient Boosting Machines (GBM) overcome this shortcoming by combining numerous decision trees in order to achieve generalization and nonlinearity \autocite{breiman_random_2001}. Retail demand forecasting literature indicates that RF models are better than the conventional regression models in their accuracy and robustness, especially when high-dimensional noisy data are to be considered \autocite{mitra_comparative_2022}.

Another popular method that is applied in classification and regression tasks is Support Vector Machines (SVMs). They are useful in the identification of nonlinear patterns through the use of kernel functions. SVMs may categorize the suppliers or the shipments as a high-risk or a low-risk group, according to the multidimensional variables like geopolitical indicators, shipment delays, and financial variables \autocite{zhao_study_2021}. SVMs however necessitate a lot of hyperparameter optimization and can be computationally inefficient with large volumes of data.

Recurrent Neural Networks (RNNs) and Long Short-Term Memory (LSTM) networks of Artificial Neural Networks (ANNs) and Deep Learning (DL) have become the focus of time-series prediction and dynamic inventory management. The RNNs are able to capture sequential dependencies and can therefore be used to predict consumer demand patterns or shipment delays. LSTMs specifically are able to retain long-term dependency as well as cope with anomalous changes, which provides high predictive rates of demand forecasting in FMCG (Fast-Moving Consumer Goods) markets \autocite{yu2023}. Nevertheless, they consume high computational resources and large volumes of labelled data making them less useful in small-scale operations \autocite{NA2024313, LIU2024}.

Recent research also notes that Reinforcement Learning (RL) can be implemented in the supply chain to make adaptive decisions. The RL agents are taught the best actions (e.g. replenishment quantity, transport scheduling) using trial and feedback, as a simulation of dynamic and uncertain environments. Some uses are optimization of routes and inventory allocation, in which RL continuously improves methods to reduce costs and delays \autocite{chowdhury2025systematic}. Still, inferential quality and time to convergence are the most important issues with RL models.


\subsection{Hybrid and Ensemble Methods}

Hybrid analytics techniques involve the use of more than one algorithm or machine learning with established optimization and simulation techniques to ensure higher accuracy and flexibility. Ensemble methods like Bagging, Boosting and Stacking combine the result of multiple base models in order to generate a more robust meta-model. An example is that Extreme Gradient Boosting (XGBoost) and Light Gradient Boosting machine (LightGBM) have been shown to have a higher predictive accuracy in terms of sales and inventory than single models \autocite{ke2017lightgbm}.

The second way of hybridity is the use of AI models together with Operations Research (OR) optimization tools. Indicatively, predictive models have the capability to determine the future demand and optimization tools such as Mixed Integer Linear Programming (MILP) or Genetic Algorithms (GA) can select the most appropriate inventory and transport assignments to make based on the predicted demand. Such integration aids in predictive as well as prescriptive decisions in supply chains \autocite{Wamba2022AIenabled}.

The other hybrid method is a combination of simulation modelling (e.g., system dynamics, agent-based modelling) and AI algorithms. Such systems have the ability to emulate the occurrences of disruptions, like port delays or suppliers’ failure and apply AI to optimize responses in real time. As it was proven by \autocite{ivanov_digital_2025}, digital twins simulations in combination with ML predict supply chain resilience to diverse stressors. The hybrid models, however, have complexity in data integration, model transparency and cost of computation which may pose an obstacle to scalability.


\section{Review of Related Studies and Applications}


The application of Artificial Intelligence (AI) and analytics in supply chain management has been much discussed within the various sectors and especially consumer goods and logistics. In this section, the reviews are critical of key empirical studies, industry applications, and research trends regarding the AI-driven and hybrid analytics systems in global supply chains. The review will focus on establishing generalized findings and issues and implications to the creation of a hybrid AI-based supply chain analytics solution in the consumer goods.

\subsection{Overview of AI-Driven Supply Chain Studies}

During the last ten years, the scholarly and commercial interest in AI-enabled supply chain analytics has increased exponentially. The studies show that AI enhances visibility of supply chains, prediction, and resiliency. In one example, \autocite{ivanov_digital_2025} have revealed that machine learning and predictive analytics enable minimizing uncertainty by enhancing the global logistics system responsiveness. Equally, \autocite{Wamba2022AIenabled} demonstrated that AI leads to a more efficient decision-making process by converting descriptive information into a course of action.

The implementation of the Big Data Analytics (BDA) and predictive modelling can also be supported empirically as enablers of supply chain agility. \autocite{kache2017challenges} argue that the data integration among various nodes (suppliers, warehouses, and customers) will help in real-time synchronization. This enables companies like Walmart and Procter and Gamble to lower lead times and enhance supply demand equilibrium. Based on case studies, companies that implement AI forecasting models (e.g., Unilever predictive demand model) record up to 20 percentage points increase in inventory turnover and up to 30 percent decline in the amount of unused stocks \autocite{Min2022BigDataAI}.

Besides, hybrid systems combining AI with more traditional methods of optimization or simulation are more adaptable. As noted by Ivanov (2023)\autocite{IVANOV2023108938}, disruptions in physical supply chains can be simulated with the help of digital twins, which are computer-based replicas of real-life supply chains, and the corrective actions should be assessed, and this greatly enhances supply chain resilience. Due to the implementation of digital twins and predictive modelling, it is possible to conduct the analysis under the conditions of uncertainty in the form of a “what-if analysis with the assistance of which the managers can predict how geopolitical instability or disruption caused by the pandemic will affect product flow.


\subsection{Studies on Data Integration and Preprocessing}

Studies note that data diversity is one of the critical hurdles of implementation AI to supply chains. Research by \autocite{Jahin_2024} emphasizes that predictive performance directly relies on data quality. They found preprocessing tools such as feature extraction, normalization and dimensionality reduction to be essential in model accuracy. Unstructured data in the form of sensor readings, reviews of texts, and social signals, among others, has to be transformed into structured forms in AI-based supply chain forecasting.

The concept of data fusion, referring to the combination of real-time IoT data with ERP and sales systems, has become an essential approach to the attainment of end-to-end visibility in the consumer goods logistics. According to the Future of SCM report, blockchain-based data structures are implemented by other large corporations like IBM and Maersk to create links between shipping documents, container tracking, and custom approvals to minimize delays, and fraud \autocite{fi11070161}. In a similar manner, the partnership between Walmart and IBM in blockchain traceability enhanced food safety since it cut down on the time food products are tracked by seven days to only 2.2 seconds. These practical examples are confirmation that data integration and secure transparency is not just a technical necessity but also a strategic facilitator of trust in consumers and regulatory compliance.


\subsection{Machine Learning and Optimization Applications}

The application of machine learning techniques has been widespread to the important supply chain operations like demand forecasting, procurement and optimization of transportation. A research conducted by \autocite{thejovathi_evaluating_nodate, kang_sales_2023} revealed that having nonlinear relationship capabilities, Random Forest (RF) and Gradient Boosting algorithms are more efficient in predicting the sales of FMCG than linear regression. Deep learning networks like Long Short-Term Memory (LM) networks have been effectively used to forecast short term demand changes and inventory. \autocite{yu2023} documented that LSTM-based prediction increased the accuracy of the prediction by 15-25 percent in comparison with ARIMA prediction models in fast-moving consumer goods.

Adaptive logistics studies have been studied with a focus on optimization via Reinforcement Learning (RL). It allows systems to learn the best decisions with simulation and feedback that is made possible by RL. The study by \autocite{chowdhury2025systematic} showed the ability of RL to enhance routing and fleet usage and reduce the transport expenses by up to 18 percent in experimental logistics. Nevertheless, it is not widely used in practice because of the high cost of computations and constant retraining.


\section{Identified Research Gaps and Justification}

Even though recent research proves that recent scholarship makes a considerable leap in implementing AI and data-driven analytics towards supply chain management, there are still multiple important gaps that support the rationale behind the current research. To begin with, the bulk of existing research has been done on individual-functional or limited uses of AI, e.g., demand forecasting, routing optimization, or inventory control, and not a unified, hybrid framework that brings together predictive, prescriptive, and human-validated analytics throughout the entire supply-chain lifecycle. Both \autocite{ivanov_digital_2025} and \autocite{Wamba2022AIenabled} admit that the following step of development is the cross-functional integration, but there is not much empirical evidence on how hybrid architectures can be implemented in consumer-goods settings.

Second, there is a gap in data-governance and interoperability, which the literature shows has persisted. Researchers like \autocite{kache2017challenges} can point to data fragmentation, the lack of standards, and weak transparency as the factors that diminish the reliability of algorithms. Although blockchain and IoT-based traceability systems demonstrated their potential, there is a lack of studies on how the technologies can be integrated with AI analytics to formulations of unified and interoperable ecosystems enabling trust, data protection, and ethical standards.

Third, the human-AI co-operation aspect, which is the focus of hybrid intelligence, has not been well investigated in consumer-goods supply chains. Current models tend to presuppose complete automation, neglecting the socio-technical interaction that is necessary to make contextual decisions, engage in ethical monitoring and explain them \autocite{arrieta2020explainable}. This creates a gap in knowledge as how to incorporate human judgment in AI-influenced decisions to find the balance between efficiency and accountability.

Fourth, there is a lack of theorisation of a sustainability connection between the AI-enabled analytics and the environmental or social context. Majority of AI research is about efficiency or cost reduction, but few evaluate the way predictive analytics can actually assist in ESG measures like carbon-footprint minimisation, reduction of waste, or responsible sourcing \autocite{nogueira2021sustainable}. It is right to fill this gap in line with the increasing international focus on green and ethical supply-chain practices.

Lastly, there exists a contextual gap in the manifestation of emerging-market supply chain. North American or European contexts are the basis of much of the empirical evidence, and there is not much understanding of how hybrid AI systems might be scaled to environments with limited data and resources like infrastructure, which are common in consumer-goods markets in Asia-Pacific.

This paper bridges these gaps by proposing a hybrid AI-based supply-chain analytics system with predictive and prescriptive intelligence, transparency with blockchain capabilities, and the human-in-the-loop decision-making process. The new knowledge brought by the proposed model is that it would show how data science, sustainability metrics, and managerial judgment can be integrated into a comprehensive, ethically responsible analytical ecosystem. By so doing, it fosters theoretical knowledge as well as practical implementation measures of next-generation supply chain of the consumer-goods industry.


\section{Chapter Summary}

The chapter has integrated a large body of both theoretical and empirical research relating to AI-driven and hybrid supply-chain analytics. It laid the formal ground with Resource-Based View, Dynamic Capabilities, and Socio-technical Systems theories, and examined a very broad spectrum of analysis mechanisms - including machine learning, deep learning, optimization and hybrid combinations. Related studies and industry applications were also discussed in the chapter, with successes in accuracy of the forecast, data integration, and digital-twin simulation and lack of success through issues of data quality, interpretability, and ethical governance.

Based on these observations, a number of gaps in research were identified with regards to integration, interoperability, human collaboration, sustainability, and applicability in emerging markets. The results were used to construct the hybrid AI-oriented framework used in the current study, which intends to provide consumer-goods supply-chain with an adaptable, transparent, and sustainable analytics solution. The following chapter will hence map the research methodology, covering the research design, methods of data-collection, analysis procedures, and strategies of validation employed in order to test empirically this proposed model.